% Proposed B.102: actual cover comparison, with no obstruction vanishing claim.
\statement{B.102}{Proved}{The actual nodal cover realizes the full syzygy module and identifies its residual map with a Smith--Tate edge map}
\label{stmt:B102}
Retain the actual symmetric complement, based involution, and discrete
modules of B.89--B.92. Let $C$ be the pullback of $Z$ to the universal
cover of the both-family filled surface $X$. Thus
\[
 \pi_1(C)=K_s,\qquad M=N_A\cap N_B=[N_A,N_B],\qquad
 E=K_s/M=R_A\times R_B.
\]
Let $C_E\to C$ be the actual connected cover associated with $M$,
and use the based lift of $F$ fixing the selected source basepoint.
There is an $E$-equivariant and $F$-equivariant nodal CW model of this
cover whose cellular complex is
\begin{equation}\label{eq:B102-cover-complex}
 P_{\mathrm{node}}\xrightarrow{\,\partial_2d\,}
 C_1(T_A\times T_B)\xrightarrow{\,\partial_1\,}
 C_0(T_A\times T_B).
\end{equation}
Here $T_A,T_B$ are the actual paired Cayley trees of B.91, and
$\partial_2:P=C_2(T_A\times T_B)\to C_1(T_A\times T_B)$ is
injective. The product-tree squares themselves are not attached;
their chains encode the boundaries of the actual transported node
cells. In particular, with the full coefficient actions retained,
\begin{equation}\label{eq:B102-cover-homology}
 H_1(C_E,\mathbb Z)=M^{\mathrm{ab}}=P/\overline L,
 \qquad H_2(C_E,\mathbb Z)=Z_2=\ker d,
 \qquad \pi_0(C_E^F)=R/I.
\end{equation}
The last set has the actual diagonal left $R$-action; $I\subseteq R$
is the source subgroup of B.92.

Place $P_{\mathrm{node}}$ in cochain degree $-2$ and $P$ in degree
$-1$, and write $D=[P_{\mathrm{node}}\xrightarrow d P]$.
Integral Tate hypercohomology gives
\begin{equation}\label{eq:B102-hyper-identification}
 \widehat{\mathbb H}^{0}(C_2,D)
   \cong I_{\mathbb F_2[R]}/J_I
   \cong \widetilde H_0(C_E^F,\mathbb F_2),
 \qquad J_I=\sum_{i\in I}\mathbb F_2[R](i-1).
\end{equation}
Under these identifications, B.92's actual map $\kappa$ is the first
map in the natural exact sequence
\begin{equation}\label{eq:B102-edge-sequence}
 \begin{aligned}
 H^2(C_2,Z_2)&\xrightarrow{\ \kappa\ }
       I_{\mathbb F_2[R]}/J_I
       \longrightarrow H^1(C_2,M^{\mathrm{ab}})\\
 &\longrightarrow H^3(C_2,Z_2)
       \longrightarrow\widehat{\mathbb H}^{1}(C_2,D).
 \end{aligned}
\end{equation}
The middle map sends $r^{-1}I-I$ to the right-defect class for the
diagonal target $(r,F(r))$. Thus its kernel is exactly
$\operatorname{im}\kappa$, with no injectivity assertion.

If $f:S^2\to C_E$ is an actual equivariant map for the orientation
preserving half-turn of $S^2$, with fixed poles $n,s$, then
\begin{equation}\label{eq:B102-sphere-value}
 \kappa\bigl([f_*[S^2]]\bigr)
   =[\operatorname{comp}(f(n))]-[\operatorname{comp}(f(s))]
       \quad\hbox{in }\widetilde H_0(C_E^F,\mathbb F_2).
\end{equation}
Such a sphere with poles in different components exists if and only
if there is an actual $F$-fixed element of $K_s$ with diagonal image
$(r,F(r))$ for some $r\notin I$. This is an equivalence, not an
existence claim. Invariant ordinary homology classes are not silently
replaced by equivariant sphere maps.

Finally, the spherical image $S_{\mathrm{sph}}$ fits into the actual
equivariant Hopf sequence
\begin{equation}\label{eq:B102-Hopf}
 \pi_2(C_E)\longrightarrow Z_2
       \longrightarrow H_2(M,\mathbb Z)\longrightarrow0,
 \qquad S_{\mathrm{sph}}=
       \operatorname{im}\bigl(\pi_2(C_E)\to Z_2\bigr).
\end{equation}
The nonspherical quotient $H_2(M,\mathbb Z)$ remains part of the
obstruction problem. No vanishing of $\kappa$, global asphericity,
nonlinear fixed lifting, or full profinite coefficient comparison is
asserted.

\begin{proof}
\emph{B.102.1. What the ordinary Brown--Loday comparison retains.}
For this paragraph denote the two one-sided open pullbacks by
$U=\widetilde Y$ and $V=\widetilde{F(Y)}$. Their intersection is $C$,
and their union is
\[
 W=U\cup V=\widetilde X\setminus\{\text{lifted mixed nodes}\}.
\]
Both $U$ and $V$ are $K(\text{free},1)$ spaces by the actual bundle
description of B.89. The filled universal cover $\widetilde X$ is
contractible. Its deleted nodes form a locally finite discrete set
in a real four-manifold. General position for compact maps of
dimension at most three shows that $W$ is two-connected. Excision,
followed by Hurewicz, gives
\begin{equation}\label{eq:B102-node-spheres}
 \pi_3(W)=H_3(W,\mathbb Z)
        =\mathbb Z[\{\text{lifted nodes}\}].
\end{equation}
The open cover is a homotopy pushout after mapping-cylinder
replacement. The homotopy-fiber theorem and its fibration sequences
give
\begin{equation}\label{eq:B102-BL-sequence}
 \pi_2(C)\longrightarrow\pi_3(W)
  \longrightarrow N_A\otimes N_B
  \xrightarrow{[\ ,\ ]}M\longrightarrow1.
\end{equation}
The tensor product is the nonabelian tensor product with the actual
mutual conjugation actions. The last image is $M$ by B.89.
The needed homotopy-fiber calculation is
\citet[Theorem~3.1 and Corollaries~3.2--3.4, pp.~318--319]{BBrownLoday1987};
it uses $\pi_2(U)=\pi_3(U)=\pi_2(V)=\pi_3(V)=0$, and does not
assume $\pi_2(C)=0$.

Replace $C$ by $K(K_s,1)$ and form the classifying-space homotopy
pushout $W_{\mathrm{cl}}$. The same theorem and B.89 give
\[
 \pi_1(W_{\mathrm{cl}})=1,\qquad
 \pi_2(W_{\mathrm{cl}})=M/[N_A,N_B]=0,\qquad
 \pi_3(W_{\mathrm{cl}})=\ker(N_A\otimes N_B\to M).
\]
Mayer--Vietoris identifies $H_3(W_{\mathrm{cl}})$ with $H_2(K_s)$,
because its two other vertices are graphs up to homotopy.
Hurewicz is an isomorphism in degree three. Naturality therefore
identifies the tensor kernel with $H_2(K_s)$ and the comparison
\eqref{eq:B102-BL-sequence} with the ordinary Hopf sequence
\begin{equation}\label{eq:B102-ordinary-Hopf}
 \pi_2(C)\longrightarrow H_2(C)
     =\mathbb Z[\{\text{lifted nodes}\}]
     \longrightarrow H_2(K_s)\longrightarrow0.
\end{equation}
Each of $N_A,N_B$ acts trivially on the tensor kernel by
\citet[Proposition~2.3(e), pp.~314--315]{BBrownLoday1987}.
They generate $K_s$, so this kernel has trivial $K_s$-action.
It cannot retain the full $E$-module $Z_2$ of the statement.

The signs in this distinction matter. Holomorphic interchange of
the two branches preserves the real orientation of a local
four-ball and its linking three-sphere. It therefore acts on
\eqref{eq:B102-node-spheres} by positive permutation. Interchanging
the two open sets in Mayer--Vietoris introduces a minus sign in
the connecting map. Equivalently, interchange of the two circle
factors reverses the orientation of the local two-torus. Thus the
node basis of $H_2(C)$ has negative permutation action. The tensor
action compatible with the commutator map is
\[
             F_T(m\otimes n)=(F(n)\otimes F(m))^{-1}.
\]
The tensor relations verify this formula and its compatibility
with $F([m,n])$. Consequently the linking-sphere map to this
tensor kernel is anti-equivariant under branch interchange,
or equivariant after the sign twist. No positive node-sphere
action is substituted for the action on commutator two-cells.

\emph{B.102.2. The actual relative homotopy model at a node.}
Work first over a sufficiently small disk about one of the actual
critical values in the original finite puncture family. No retained
punctures collide there. Choose disjoint moving fiber disks about
the colliding pairs, avoiding every other trajectory. The exterior
punctures and disk boundaries trivialize over the entire base disk
by the finite-family isotopy of B.91. The exterior is a product
with a punctured planar surface; its new boundary circles are
primitive members of a free peripheral basis.

Inside a collision disk the two transverse branches have holomorphic
coordinates $x,y$ with equations $x=0$ and $y=0$. A small round
normal neighborhood minus the branches retracts radially to the
Hopf-link complement. Its angular coordinates and radial ratio
give the usual homotopy equivalence with $T^2$, hence a
$K(\mathbb Z^2,1)$. Concentric node boxes and annular collars give
the same pair. The collar generator is the actual product $ab$
of the two positive meridians and injects in this $\mathbb Z^2$.

The local complement is consequently a tree of aspherical spaces
with injective edge maps. Its universal cover is a tree of
contractible vertex and edge spaces, hence is contractible.
Eliminating the primitive exterior boundary generators gives
\begin{equation}\label{eq:B102-local-group}
 \pi_1(\text{local block})
      =F_{\mathrm{free}}*\mathop{*}_{\nu\text{ in the block}}\mathbb Z^2.
\end{equation}
Replace the generic fiber by its actual peripheral graph and attach
one two-cell along each transported commutator $[a,b]$. In a local
adapted marking this is a wedge of the same free graph and one
torus per collision. It is aspherical with the group
\eqref{eq:B102-local-group}. Map its generic graph to the actual
nearby fiber and fill its commutator cells by the local torus maps.
This prescribed map of pairs induces the identity on the displayed
fundamental group, so it is a homotopy equivalence. Its restriction
to the generic graph is the specified fiber inclusion; the graph
has no higher homotopy obstruction to the required based pair
homotopy. This verifies the relative attaching comparison, rather
than deducing a homotopy type from van Kampen alone.

Pull this finite relative equivalence through the actual retained
$F_5$ cover. Covering homotopy lifting preserves it as a relative
equivalence. The generic graph is the actual generic-cover kernel
graph, and all retained $F_5$ translates of the local node cells
are included. In particular repeated meetings of lifted components
retain their separate cells and their actual based words.

\emph{B.102.3. Equivariance and the global relative construction.}
B.92 normalizes the regular-fiber involution to $z\mapsto-z$
and constructs its actual equivariant regular-path transports.
At a fixed node, choose one branch coordinate $x$ and its $F$-image
$y$. They form a holomorphic coordinate system because the branches
are transverse, and $F$ exchanges $x,y$. Choose the torus model
symmetrically. With the actual nearby source basepoint and paired
meridian arcs, its fixed diagonal circle is the source peripheral
$ab$. The fixed subspace of the node complement is the punctured
diagonal disk; this circle map is a homotopy equivalence on that
fixed subspace as well as on the underlying block.

In the selected universal old-fiber cover each regular fiber has
exactly one fixed point. Thus over a critical value there is at
most one fixed lifted node; the others occur in pairs. Without a
fixed node the local fixed subspace is the unpunctured source disk
and its model is contractible. With a fixed node both fixed
subspaces have the circle homotopy type just described. Choose
one local filling for each nonfixed pair and use its actual
$F$-image for the other. Equivariant exterior isotopies retain the
source whisker. The generic paired graph comes from B.91's quotient
orbifold arcs and has exactly the selected fixed vertex.

These local pair maps are equivalences both on underlying spaces
and on fixed subspaces, compatibly with the common generic fiber.
After equivariant mapping-cylinder replacement they are equivalences
of equivariant CW pairs. For completeness, the equivariant Whitehead
step here is the cellular induction: a cell has stabilizer $1$ or
$C_2$, and extension of a homotopy inverse and its homotopies uses
the vanishing relative homotopy groups of the corresponding fixed
subspaces. Underlying fundamental-group agreement alone would not
supply this step.

Use B.92's common locally finite tree of critical-value paths,
retaining its source peripheral paths. Regular strips use the
equivariant finite-family trivializations, and critical disks use
the relative maps above. Mapping cylinders make the gluing
cofibrant. Each compact stage preserves both underlying and fixed
homotopy equivalences. Exhaust the universal old base by compatible
such stages. Compact loops, spheres and homotopies meet only
finitely many lifted charts. The union is a weak equivalence on
underlying and fixed subspaces and hence an equivariant CW homotopy
equivalence. This produces a based nodal two-complex $V_{\mathrm{nod}}$
for the actual $C$, with its actual paired generic graph and every
transported commutator cell. Its fixed graph models the chosen
punctured source. No global asphericity is used; B.100 is compatible
with the nonzero second homotopy of this two-complex.

\emph{B.102.4. Lifting the comparison to the actual coefficient cover.}
The based comparison induces the actual quotient $K_s\to E$.
Lift it through that quotient. Uniqueness of based lifting makes
the lift commute with deck transformations and with $F$. The
cover of the paired generic graph is
$(T_A\times T_B)^{(1)}$. The actual node cell for $\nu$ has edge
boundary $\partial_2z_\nu$ in this graph, by the defining full
Fox-chain comparison of B.91. Therefore its lifted chain complex
is precisely \eqref{eq:B102-cover-complex}. Contractibility of the
product of the two trees gives injectivity of $\partial_2$ and
exactness of its augmented cellular complex. This proves the
first two identifications in \eqref{eq:B102-cover-homology}, with
both coefficient actions. The cover has not been extended over
$W$ or either filling; its nontrivial $E$-monodromy is retained.

By B.87 and B.92 the selected fixed locus in $C$ is the punctured
chosen source $\widetilde S$. At a fixed basepoint its fixed fiber
labels in $C_E$ are exactly the diagonal elements $(r,F(r))$.
Lifting source paths identifies these labels by the diagonal image
of $\pi_1(\widetilde S)$, which is the actual subgroup $I$.
Each component containing one fixed lift is fixed pointwise by
uniqueness of lifting the identity on $\widetilde S$, and every
fixed point projects into $\widetilde S$. Hence
\begin{equation}\label{eq:B102-fixed-components}
 \pi_0(C_E^F)=R/I,\qquad H_0(C_E^F,\mathbb Z)=\mathbb Z[R/I].
\end{equation}
Each component is a covering of a punctured disk. Its first
homology is free abelian and it has no higher homology. This
argument does not assume that the source fundamental group
injects in $K_s$. The source graph comparison lifts on every
fixed component of the nodal model.

\emph{B.102.5. Tate totalization and the literal residual map.}
Use cochain degree equal to minus cellular dimension. In the
reduced version of \eqref{eq:B102-cover-complex}, the last term is
the augmentation kernel of $C_0(T_A\times T_B)$. The map
$\partial_2$ gives a quasi-isomorphism from $D$ to this reduced
complex: its quotient has just two terms, identified by the
isomorphism from $C_1/P$ to that augmentation kernel. All maps
commute with $E$ and $F$.

Apply the complete two-periodic $C_2$ cochain resolution, with
alternating maps $1-F$ and $1+F$, and totalize. The cellular
direction is bounded, so each total degree uses finitely many
cellular degrees. The row and column filtrations converge, even
for these infinite-rank discrete modules. The same argument shows
that bounded equivariant quasi-isomorphisms induce isomorphisms on
this Tate hypercohomology; no completed support convention is used.

Both $P$ and $P_{\mathrm{node}}$ have zero even Tate cohomology.
Their odd cohomology is the binary module on the fixed indices
of their negative permutation actions. B.91--B.92 identify the
image between those odd groups as $J_I$. The two-term totalization
therefore gives the first isomorphism of
\eqref{eq:B102-hyper-identification}. Its homology truncation has
$H^{-2}(D)=Z_2$ and $H^{-1}(D)=M^{\mathrm{ab}}$. The exact triangle
of these homology objects gives \eqref{eq:B102-edge-sequence}.

To identify its first map without a convention ambiguity, represent
a class of $H^2(C_2,Z_2)$ by $z\in Z_2$ with $Fz=z$, modulo norms.
Since $H^2(C_2,P_{\mathrm{node}})=0$, write $z=(1+F)p$ for a node
chain $p$. Then $d(p)$ is anti-invariant, and the edge image is
its diagonal Tate class modulo $J_I$. This is exactly the
representative formula for $\kappa$ in B.92. Changing $p$ adds an
anti-invariant node-chain image, hence changes the value only by
$J_I$; changing $z$ by a norm in $Z_2$ changes no value. Thus the
comparison uses the actual residual map, not a new map with only
the same source and target.

There is also a direct Smith-chain interpretation. Subdivide every
fixed commutator disk along its reflection diagonal. Its two open
halves are exchanged and its diagonal is pointwise fixed. Subdivide
paired cells compatibly. Cellular chains relative to the fixed
graph are then sums of free $\mathbb Z[C_2]$-modules. Their Tate
hypercohomology is zero, because this relative complex is bounded
and every such module is Tate-acyclic. Inclusion of the fixed
graph is consequently an isomorphism on reduced Tate
hypercohomology. The graph has trivial action and free integral
homology in degrees zero and one. Its degree-zero reduced Tate
hypercohomology is therefore
\[
 \widetilde H_0(C_E^F,\mathbb F_2)
   =\ker\bigl(\mathbb F_2[R/I]\to\mathbb F_2\bigr)
   =I_{\mathbb F_2[R]}/J_I.
\]
This proves the second isomorphism of
\eqref{eq:B102-hyper-identification}.

The endpoint and defect conventions are retained. A path from
vertex $1$ to vertex $(r,F(r))$ gives the left defect
$F(x)x^{-1}$ and the coordinate $r-1$. The right defect
$x^{-1}F(x)$ is its diagonal $r^{-1}$ translate and has coordinate
$r^{-1}-1$ over $\mathbb F_2$, as in B.91. Thus the map in
\eqref{eq:B102-edge-sequence} evaluated at $r^{-1}I-I$ gives the
right-defect class of target $(r,F(r))$. It does not relabel the
endpoint $rI$ by $r^{-1}I$ or alter the diagonal left action.

\emph{B.102.6. What the fixed-graph comparison leaves open.}
The fixed graph has dimension one, but this does not kill the
first map of \eqref{eq:B102-edge-sequence}. An invariant two-cycle
can contribute through $H^2(C_2,H_2(C_E))$. A submodule of a
Tate-acyclic chain module need not be Tate-acyclic. In particular,
the absence of pointwise fixed two-cells after subdivision is not
a vanishing proof for $\kappa$.

B.100 supplies an actual nonzero sphere and rules out a proof
based on contractibility of the universal cover. That fact alone
does not give a nonzero value of the present map. The ordinary
comparison in B.102.1 forgets the coefficients that the actual
cover retains. The problem remains to evaluate
\eqref{eq:B102-edge-sequence}, or to exclude differences of
distinct coset basis vectors from its image. Such an exclusion
would still require the separate nonlinear and compact
fixed-lifting comparisons to address those stronger conclusions.

\emph{B.102.7. Equivariant spheres and the nonlinear projected test.}
Let $C_2$ act on $S^2$ by a half-turn about its vertical axis.
It preserves the fundamental class and fixes exactly its two
poles. Its reduced Smith--Tate edge is the isomorphism
\[
 H^2(C_2,H_2(S^2,\mathbb Z))=\mathbb F_2
       \longrightarrow\widetilde H_0(\{n,s\},\mathbb F_2)
       =\mathbb F_2,
\]
taking the fundamental class to $n-s$. This follows directly by
dividing the sphere into two exchanged lunes, or by the same
bounded totalization, since $H_1(S^2)=0$. Naturality proves
\eqref{eq:B102-sphere-value}. For poles over vertices $1$ and
$(r,F(r))$, its value is $rI-I$. Translating the entire sphere by
the $F$-fixed deck element $(r^{-1},F(r^{-1}))$ gives instead
$r^{-1}I-I$, matching the stated right-defect convention. These
two values differ by that specified invertible left translation;
they are not identified without it.

For the nonlinear interpretation choose fixed paths within the
two pole components to points over the original source basepoint.
After diagonal translation, their labels are $1$ and $(r,F(r))$.
A meridian arc on the sphere and its $F$-image are homotopic
relative endpoints. Their images in $C_E$ project to the same
based element of $K_s$, so that element is literally $F$-fixed
and has the required diagonal head. Distinct pole components
mean $r\notin I$.

Conversely, let $g\in K_s$ be literally $F$-fixed with diagonal
image $(r,F(r))$. Lift a representative loop to a path $\gamma$
in $C_E$ from the chosen basepoint to its diagonal translate.
The equality $F(g)=g$ supplies a based homotopy between
$\gamma$ and $F\gamma$, lifting relative endpoints. Use this
homotopy on one lune and its $F$-image on the other. Their
boundaries agree and give an actual equivariant half-turn sphere
with the asserted poles. This proves the equivalence in the
statement. General projected heads in $H_b$ reduce to this
zero-head case by multiplying by an actual fixed source lift,
as in B.92. An arbitrary invariant class in $Z_2$ need not be
represented by such an equivariant sphere; the equivalence
does not supply a counterexample or a vanishing theorem.

\emph{B.102.8. The separate nonspherical quotient.}
The ordinary Hopf sequence for the actual connected cover $C_E$
is natural for $F$ and gives \eqref{eq:B102-Hopf}. Covering
invariance identifies its second homotopy group with $\pi_2(Z)$,
but its Hurewicz image need not be all of $Z_2$ and need not
identify with $\pi_2(Z)$ injectively. Consequently even all deck
and involution translates of known sphere families do not
automatically determine the entire domain of $\kappa$.

For example, a valid sufficient two-part condition for vanishing
would be
\begin{equation}\label{eq:B102-conditional-spherical}
 \begin{gathered}
 \kappa\circ\bigl(H^2(C_2,S_{\mathrm{sph}})
                      \to H^2(C_2,Z_2)\bigr)=0,\\
 H^2\bigl(C_2,H_2(M,\mathbb Z)\bigr)=0.
 \end{gathered}
\end{equation}
Indeed the long exact sequence of
$0\to S_{\mathrm{sph}}\to Z_2\to H_2(M,\mathbb Z)\to0$
would make the first displayed cohomology map surjective.
Neither premise of \eqref{eq:B102-conditional-spherical} is
established. This keeps the nonspherical group-homology quotient
explicit rather than replacing every two-cycle by a transported
spherical identity.

\begingroup\raggedright
The complete actual-space and sign comparison is retained in
\path{work_A/actual_nodal_cover_Smith_comparison_19h.md}.
The primary Brown--Loday theorem, tensor action and classifying-space
comparison were independently reread; the relative equivariant
geometric model and the Tate truncation were reviewed separately.
These are discrete full-coefficient constructions. No finite
truncation of the old-base orbit is substituted for their
actual node module.\par
\endgroup
\end{proof}
